\documentclass{article}
\usepackage{graphicx}
\usepackage{float}
\usepackage[a4paper, margin=1in]{geometry}
\usepackage[style=ieee,backend=bibtex]{biblatex}
\usepackage{setspace}
\usepackage{hyperref}
\usepackage{spreadtab}
\usepackage{amsmath}

\newfloat{eqn}{thp}{lop}

\hypersetup{
    colorlinks = true,
    urlcolor = blue,
    linkcolor = black,
    citecolor = black
}

\onehalfspacing

\addbibresource{references.bib}



\begin{document}

Specifications for a few common commercial drones were compiled to get an overview of typical payload capacity and performance.
Most are consumer drones often used for photography, a few higher end drones used for specialty applications like search and rescue.
Three specifications were collected for each drone, if available:
\begin{itemize}
    \item Weight, listed in table \ref{tab:weights}.
    \item Maximum payload, listed in table \ref{tab:payload}.
    \item Maximum wind tolerance, listed in table \ref{tab:wind}.
\end{itemize}

The weights show that most low grad drones weight under 1 kg. The special purpose drones can approach 4 kg, such heavy designs will not be targeted for this
thesis.

The rated payload was calculated as the maximum total takeoff weight minus the drones weight. The DJI Mavic products had no maximum takeoff weight aside from 
the products weight, so no maximum payload is listed for them. 

The wind resistance values are typically around 10 m/s for low grad drones, increasing to 12 to 14 for the heavier models. The DJI Mavic Pro and Mavic 3 Pro
has no wind resistance listed in their specifications, a product FAQ stated that they are both capable of operating in "Level 5" wind, which is 8 to 11 m/s.
For the proposed device to be useful, it has to hold a drone in at least as much wind as the drone can fly in. Ideally it would support use in higher wind speeds, 
so the drone could perch as a means to operate through harsh conditions. The target of this thesis will be to double the wind resistance of a small drone when perched,
for a wind resistance of 20 m/s.

The force that wind will exert on the device is highly dependent on the specific drones drag coefficient. This must be found for each design experimentally or in
simulation, which is too time consuming of a process for this thesis. Experimental data for one drone design by Hattenberger et al.
found drag coefficients between 0.14 and 0.24 \cite{hattenberger_evaluation_2023}. This was for a drone in flight rather than stationary, but for this 
concept development 0.24 is assumed to be a useful estimate of a worst case drag coefficient for a typical drone. With a maximum wind speed of 20 m/s, 
the force this device must hold (in addition to the drones weight) is 4.8 N, equivalent to an additional 490 g.

The low end of payload capacity listed is 130 g. For this device to be used on such a drone, while leaving 70\% of the rated payload to be used for 
sensors required by the application, the device must weigh less than 39 grams. With a 1 kg drone and additional sensor payload, the device must be capable of supporting
at least 1620 g.


\begin{table}[H]
    \centering
    \caption{Specified Weight}
    \begin{tabular}{| c | c |}
        \hline
        {Make \& Model} & Weight (g) \\
        \hline
        DJI Mavic Pro & 734 \\
        DJI Mavic 3 Pro & 958 \\
        DJI Mavic 3 Enterprise & 920 \\
        DJI Mavic Mini & 249 \\
        Parrot ANAFI USA & 500 \\
        FLIR SIRAS & 3100 \\
        Autel EVO Max 4T & 1640 \\
        DJI Matrice 30 & 3770 \\
        \hline
    \end{tabular}
    \label{tab:weights}
\end{table}


\begin{table}[H]
    \centering
    \caption{Rated Payload}
    \begin{tabular}{| c | c |}
        \hline
        {Make \& Model} & Payload (g) \\
        \hline
        DJI Mavic Pro & - \\
        DJI Mavic 3 Pro & - \\
        DJI Mavic 3 Enterprise & 130 \\
        DJI Mavic Mini & - \\
        Parrot ANAFI USA & 144 \\
        FLIR SIRAS & 3100 \\
        Autel EVO Max 4T & 379 \\
        DJI Matrice 30 & 299 \\
        \hline
    \end{tabular}
    \label{tab:payload}
\end{table}


\begin{table}[H]
    \centering
    \caption{Rated Wind Resistance}
    \begin{tabular}{| c | c |}
        \hline
        {Make \& Model} & Max Wind Speed (m/s) \\
        \hline
        DJI Mavic Pro & 10 \\
        DJI Mavic 3 Pro & 10 \\
        DJI Mavic 3 Enterprise & 12 \\
        DJI Mavic Mini & 8 \\
        Parrot ANAFI USA & 14 \\
        FLIR SIRAS & 10 \\
        Autel EVO Max 4T & 12 \\
        DJI Matrice 30 & 12 \\
        \hline
    \end{tabular}
    \label{tab:wind}
\end{table}



\printbibliography[heading=none]


\end{document}